\documentclass{article}
\usepackage[T1]{fontenc}
\usepackage{lmodern}
\usepackage{graphicx}
\usepackage{amsmath,amssymb}
\usepackage[a4paper,margin=2cm]{geometry}
\usepackage[absolute,overlay]{textpos}
\setlength{\TPHorizModule}{1mm}
\setlength{\TPVertModule}{1mm}
\usepackage[indonesian]{babel}
\usepackage{hyperref}
\setlength{\emergencystretch}{2em}
% unit: Halaman 39

\begin{document}

% === Halaman 39 (llm) ===

\section*{Ranah Steganografi}

Berdasarkan ranah operasinya, metode-metode steganografi dapat dibagi menjadi dua kelompok

\begin{itemize}
    \item \textbf{\textcolor{red}{Spatial (time) domain methods}}
    Memodifikasi langsung nilai $byte$ dari $cover-object$ (nilai $byte$ merepresentasikan intensitas/warna $pixel$ atau amplitudo)
    \par Contoh: Metode $LSB$
    
    \item \textbf{\textcolor{red}{Tranform domain methods}}
    Memodifikasi hasil transformasi sinyal dalam ranah transform (hasil transformasi dari ranah spasial ke ranah lain (misalnya ranah frekuensi).
    \par Contoh: Metode $Spread$ $Spectrum$
\end{itemize}

\end{document}
